\section{The articles of the angelic treatise}
\label{app:census}

The census follows all 118 articles in the twenty-four questions of the
\emph{Summa} treated here. Each row gives the question posed, Aquinas's
determination and the principal authority or argument of the
\emph{sed contra}. The grade belongs to the stated determination.

\Needspace{10\baselineskip}
\subsection{Question 50}
\begin{longtable}{@{}>{\raggedright\arraybackslash}p{0.31\linewidth}>{\raggedright\arraybackslash}p{0.39\linewidth}>{\raggedright\arraybackslash}p{0.23\linewidth}@{}}
\toprule
Article and question & Aquinas's determination & Grade; authority\\
\midrule\endhead
\textbf{1.} Whether an Angel Is Altogether Incorporeal? & Angels are wholly incorporeal. & Church teaching. Psalm 103:4 [Hebrew 104:4]\\
\textbf{2.} Whether an Angel Is Composed of Matter and Form? & Angels have no matter-form composition, but retain essence-existence and potency-act composition. & Thomist position. Dionysius, Divine Names IV\\
\textbf{3.} Whether the Angels Exist in Any Great Number? & Angels form an immense multitude, exceeding material multitude in Aquinas's argument. & Thomist position. Daniel 7:10\\
\textbf{4.} Whether the Angels Differ in Species? & Each angel is specifically distinct; two angels cannot be individuals of one species. & Thomist position. Aristotle, Metaphysics III, and Dionysius, Celestial Hierarchy X\\
\textbf{5.} Whether the Angels Are Incorruptible? & Angels are naturally incorruptible while remaining dependent on divine conservation. & Church teaching. Dionysius, Divine Names IV\\
\bottomrule
\end{longtable}
\Needspace{10\baselineskip}
\subsection{Question 51}
\begin{longtable}{@{}>{\raggedright\arraybackslash}p{0.31\linewidth}>{\raggedright\arraybackslash}p{0.39\linewidth}>{\raggedright\arraybackslash}p{0.23\linewidth}@{}}
\toprule
Article and question & Aquinas's determination & Grade; authority\\
\midrule\endhead
\textbf{1.} Whether the Angels Have Bodies Naturally United to Them? & Angels have no bodies naturally united to them. & Church teaching. Dionysius, Divine Names IV\\
\textbf{2.} Whether Angels Assume Bodies? & Angels sometimes assume real bodies for manifestation to human beings. & Thomist position. Augustine, City of God XVI\\
\textbf{3.} Whether the Angels Exercise Functions of Life in the Bodies Assumed? & Assumed bodies exhibit external effects of living acts, but do not perform such acts as living organisms. & Thomist position. Prior conclusion: assumed bodies lack life\\
\bottomrule
\end{longtable}
\Needspace{10\baselineskip}
\subsection{Question 52}
\begin{longtable}{@{}>{\raggedright\arraybackslash}p{0.31\linewidth}>{\raggedright\arraybackslash}p{0.39\linewidth}>{\raggedright\arraybackslash}p{0.23\linewidth}@{}}
\toprule
Article and question & Aquinas's determination & Grade; authority\\
\midrule\endhead
\textbf{1.} Whether an Angel Is in a Place? & An angel is present in a bodily place through application of its power, without dimensional containment. & Thomist position. Collect Visita, quaesumus: invocation of angels dwelling here\\
\textbf{2.} Whether an Angel Can Be in Several Places at Once? & An angel is present definitively in one operative place, which may be extended, not everywhere or in several places. & Thomist position. John Damascene, Orthodox Faith II.3\\
\textbf{3.} Whether Several Angels Can Be at the Same Time in the Same Place? & Two angels cannot be complete immediate operative causes of one and the same place in the same respect. & Thomist position. Analogy: two souls do not inform one body\\
\bottomrule
\end{longtable}
\Needspace{10\baselineskip}
\subsection{Question 53}
\begin{longtable}{@{}>{\raggedright\arraybackslash}p{0.31\linewidth}>{\raggedright\arraybackslash}p{0.39\linewidth}>{\raggedright\arraybackslash}p{0.23\linewidth}@{}}
\toprule
Article and question & Aquinas's determination & Grade; authority\\
\midrule\endhead
\textbf{1.} Whether an Angel Can Be Moved Locally? & An angel can change local presence continuously or discontinuously through successive operative contacts. & Thomist position. Analogy with Christ's soul descending to the dead\\
\textbf{2.} Whether an Angel Passes Through Intermediate Space? & Continuous motion traverses intermediate places; discontinuous angelic motion need not. & Thomist position. Argument from a process preceding its achieved change; qualified in reply\\
\textbf{3.} Whether the Movement of an Angel Is Instantaneous? & Angelic movement includes before and after, measured by continuous or discontinuous time, not one indivisible instant. & Thomist position. Argument: movement has before and after\\
\bottomrule
\end{longtable}
\Needspace{10\baselineskip}
\subsection{Question 54}
\begin{longtable}{@{}>{\raggedright\arraybackslash}p{0.31\linewidth}>{\raggedright\arraybackslash}p{0.39\linewidth}>{\raggedright\arraybackslash}p{0.23\linewidth}@{}}
\toprule
Article and question & Aquinas's determination & Grade; authority\\
\midrule\endhead
\textbf{1.} Whether an Angel's Act of Understanding Is His Substance? & A created angel's act of understanding is not its substance. & Thomist position. Only God's existence is identical with substance; operation is still more distinct\\
\textbf{2.} Whether in the Angel to Understand Is to Exist? & An angel's understanding is not its act of existence. & Thomist position. Dionysius, Divine Names IV: understanding as movement\\
\textbf{3.} Whether an Angel's Power of Intelligence Is His Essence? & An angel's intellectual power differs from its essence. & Thomist position. Dionysius, Celestial Hierarchy XI: essence, power, operation\\
\textbf{4.} Whether There Is an Active and a Passive Intellect in an Angel? & Angels lack active and possible intellect in the human, abstraction-related sense. & Thomist position. Active/possible distinction relates to phantasms, absent in angels\\
\textbf{5.} Whether There Is Only Intellectual Knowledge in the Angels? & Angelic knowledge is intellectual; no organic sense or imagination belongs to an angel. & Thomist position. Gregory the Great, Gospel homily 29\\
\bottomrule
\end{longtable}
\Needspace{10\baselineskip}
\subsection{Question 55}
\begin{longtable}{@{}>{\raggedright\arraybackslash}p{0.31\linewidth}>{\raggedright\arraybackslash}p{0.39\linewidth}>{\raggedright\arraybackslash}p{0.23\linewidth}@{}}
\toprule
Article and question & Aquinas's determination & Grade; authority\\
\midrule\endhead
\textbf{1.} Whether the Angels Know All Things by Their Substance? & Angels do not know everything through their own finite essence; other objects require species. & Thomist position. Dionysius, Divine Names IV\\
\textbf{2.} Whether the Angels Understand by Species Drawn from Things? & Angelic intelligible species are connatural gifts from God, not abstracted from objects. & Thomist position. Dionysius, Divine Names VII\\
\textbf{3.} Whether the Higher Angels Understand by More Universal Species Than the Lower Angels? & Higher angels know through fewer, more universal species, with more complete rather than vaguer cognition. & Thomist position. Dionysius, Celestial Hierarchy XII; Book of Causes\\
\bottomrule
\end{longtable}
\Needspace{10\baselineskip}
\subsection{Question 56}
\begin{longtable}{@{}>{\raggedright\arraybackslash}p{0.31\linewidth}>{\raggedright\arraybackslash}p{0.39\linewidth}>{\raggedright\arraybackslash}p{0.23\linewidth}@{}}
\toprule
Article and question & Aquinas's determination & Grade; authority\\
\midrule\endhead
\textbf{1.} Whether an Angel Knows Himself? & An angel understands itself through its own immaterial substance. & Thomist position. Augustine, Literal Meaning of Genesis II\\
\textbf{2.} Whether One Angel Knows Another? & An angel knows another by an intelligible species, distinct from the other's natural subsistence. & Thomist position. Book of Causes: intelligences know incorruptible things\\
\textbf{3.} Whether an Angel Knows God by His Own Natural Principles? & Angels naturally know God through his image in their essence, not by seeing the divine essence. & Common teaching. Romans 1:19, with argument from angelic intellectual superiority\\
\bottomrule
\end{longtable}
\Needspace{10\baselineskip}
\subsection{Question 57}
\begin{longtable}{@{}>{\raggedright\arraybackslash}p{0.31\linewidth}>{\raggedright\arraybackslash}p{0.39\linewidth}>{\raggedright\arraybackslash}p{0.23\linewidth}@{}}
\toprule
Article and question & Aquinas's determination & Grade; authority\\
\midrule\endhead
\textbf{1.} Whether the Angels Know Material Things? & Angels know material natures through connatural intelligible species. & Thomist position. A superior cognitive power includes what a lower can know\\
\textbf{2.} Whether an Angel Knows Singulars? & Angels know singulars and their individual conditions, necessary for their particular ministries. & Thomist position. Psalm 90:11 [Hebrew 91:11]\\
\textbf{3.} Whether Angels Know the Future? & Angels naturally know determined futures through causes and probable ones conjecturally; future contingents as actual require revelation. & Common teaching. Isaiah 41:23\\
\textbf{4.} Whether Angels Know Secret Thoughts? & Angels can infer thoughts from effects, but do not naturally know free thoughts and affections as interior acts. & Common teaching. Jeremiah 17:9--10\\
\textbf{5.} Whether the Angels Know the Mysteries of Grace? & Mysteries of grace are not naturally known; revelation discloses them unequally and successively. & Common teaching. Dionysius, Celestial Hierarchy VII, interpreting Isaiah 63:1\\
\bottomrule
\end{longtable}
\Needspace{10\baselineskip}
\subsection{Question 58}
\begin{longtable}{@{}>{\raggedright\arraybackslash}p{0.31\linewidth}>{\raggedright\arraybackslash}p{0.39\linewidth}>{\raggedright\arraybackslash}p{0.23\linewidth}@{}}
\toprule
Article and question & Aquinas's determination & Grade; authority\\
\midrule\endhead
\textbf{1.} Whether the Angel's Intellect Is Sometimes in Potentiality, Sometimes in Act? & Natural species are already possessed, but consideration can be successive; revelation can supply new knowledge; beatific vision remains actual. & Thomist position. Augustine, Literal Meaning of Genesis II\\
\textbf{2.} Whether an Angel Can Understand Many Things at the Same Time? & An angel simultaneously knows objects comprised under one species, not naturally all objects requiring different species. & Thomist position. Augustine, Literal Meaning of Genesis IV.32\\
\textbf{3.} Whether an Angel's Knowledge Is Discursive? & Angelic natural knowing is nondiscursive: conclusions are apprehended in principles without discovery by reasoning. & Thomist position. Dionysius, Divine Names VII\\
\textbf{4.} Whether the Angels Understand by Composing and Dividing? & Angels know propositions without acquiring knowledge by composing and dividing. & Thomist position. Dionysius, Divine Names VII\\
\textbf{5.} Whether There Can Be Falsehood in the Intellect of an Angel? & Angels do not err about naturally known essences; demons can misjudge supernatural matters through disordered application. & Thomist position. Aristotle, De anima III; Augustine, Eighty-three Questions 32\\
\textbf{6.} Whether There Is a "Morning" and an "Evening" Knowledge in the Angels? & Morning and evening name knowledge of creatures in the Word and in their own existence. & Thomist position. Augustine, Literal Meaning of Genesis IV.22,31; City of God as cited\\
\textbf{7.} Whether the Morning and Evening Knowledge Are One? & They differ in medium if evening uses connatural species; if both are through the Word they differ only by object-aspect. & Thomist position. Augustine, Literal Meaning of Genesis IV.24\\
\bottomrule
\end{longtable}
\Needspace{10\baselineskip}
\subsection{Question 59}
\begin{longtable}{@{}>{\raggedright\arraybackslash}p{0.31\linewidth}>{\raggedright\arraybackslash}p{0.39\linewidth}>{\raggedright\arraybackslash}p{0.23\linewidth}@{}}
\toprule
Article and question & Aquinas's determination & Grade; authority\\
\midrule\endhead
\textbf{1.} Whether There Is Will in the Angels? & Angels have will as intellectual appetite for good universally understood. & Church teaching. Augustine, Trinity X.11--12\\
\textbf{2.} Whether in the Angels the Will Differs from the Intellect? & Angelic will differs from both essence and intellect. & Thomist position. Will regards good; intellect also understands evil\\
\textbf{3.} Whether There Is Free-Will in the Angels? & Angels have free choice without requiring discursive deliberation. & Common teaching. Human freedom as dignity, more perfectly in angels\\
\textbf{4.} Whether There Is an Irascible and a Concupiscible Appetite in the Angels? & Angels have no concupiscible or irascible sensitive appetite; love and joy can name acts of will. & Thomist position. Aristotle, De anima III: those appetites belong to sensation\\
\bottomrule
\end{longtable}
\Needspace{10\baselineskip}
\subsection{Question 60}
\begin{longtable}{@{}>{\raggedright\arraybackslash}p{0.31\linewidth}>{\raggedright\arraybackslash}p{0.39\linewidth}>{\raggedright\arraybackslash}p{0.23\linewidth}@{}}
\toprule
Article and question & Aquinas's determination & Grade; authority\\
\midrule\endhead
\textbf{1.} Whether There Is Natural Love or Dilection in an Angel? & Angels have a natural inclination of love in their intellectual will. & Thomist position. Augustine, Trinity X.1--2; love follows knowledge\\
\textbf{2.} Whether There Is Love of Choice in the Angels? & Angels have elective love grounded in natural love of the end. & Thomist position. Merit and demerit exceed merely natural acts\\
\textbf{3.} Whether the Angel Loves Himself with Both Natural Love, and Love of Choice? & Angels love themselves naturally and by chosen goods desired for themselves. & Thomist position. Aristotle, Ethics IX.8\\
\textbf{4.} Whether an Angel Loves Another with Natural Love As He Loves Himself? & Angels naturally love one another through shared nature, similarly but not equally to self-love. & Thomist position. Ecclesiasticus 13:19\\
\textbf{5.} Whether an angel by natural love loves God more than he loves himself? & Natural love orders an angel to God as universal good above its own private good; charity remains supernatural. & Thomist position. Moral precept to love God above self belongs to natural law\\
\bottomrule
\end{longtable}
\Needspace{10\baselineskip}
\subsection{Question 61}
\begin{longtable}{@{}>{\raggedright\arraybackslash}p{0.31\linewidth}>{\raggedright\arraybackslash}p{0.39\linewidth}>{\raggedright\arraybackslash}p{0.23\linewidth}@{}}
\toprule
Article and question & Aquinas's determination & Grade; authority\\
\midrule\endhead
\textbf{1.} Whether the Angels Have a Cause of Their Existence? & Angels receive their entire existence from God. & Defined. Psalm 148:2,5\\
\textbf{2.} Whether the Angel Was Produced by God from Eternity? & Angels had a beginning and are not coeternal with God. & Defined. Proverbs 8:22\\
\textbf{3.} Whether the Angels Were Created Before the Corporeal World? & Simultaneous creation with bodies is judged more probable; prior angelic creation is not condemned. & Disputed. Genesis 1:1\\
\textbf{4.} Whether the Angels Were Created in the Empyrean Heaven? & Creation in the highest bodily heaven is fitting to angelic government, without bodily dependence. & Disputed. Strabus gloss on Genesis 1:1\\
\bottomrule
\end{longtable}
\Needspace{10\baselineskip}
\subsection{Question 62}
\begin{longtable}{@{}>{\raggedright\arraybackslash}p{0.31\linewidth}>{\raggedright\arraybackslash}p{0.39\linewidth}>{\raggedright\arraybackslash}p{0.23\linewidth}@{}}
\toprule
Article and question & Aquinas's determination & Grade; authority\\
\midrule\endhead
\textbf{1.} Whether the Angels Were Created in Beatitude? & Angels begin with natural perfection, not consummated supernatural beatitude. & Thomist position. Their possible fall excludes initial confirmation in beatitude\\
\textbf{2.} Whether an Angel Needs Grace in Order to Turn to God? & Grace is needed to turn to God as supernatural beatitude. & Common teaching. Romans 6:23\\
\textbf{3.} Whether the Angels Were Created in Grace? & Creation in sanctifying grace is judged more probable, though opinions differ. & Disputed. Augustine, City of God XII.9\\
\textbf{4.} Whether an Angel Merits His Beatitude? & Angels merited beatitude by a grace-informed act preceding enjoyment. & Thomist position. Apocalypse 21:17: measure of man and angel\\
\textbf{5.} Whether the Angel Obtained Beatitude Immediately After One Act of Merit? & Beatitude followed one meritorious act immediately in the succession of angelic acts. & Thomist position. Analogy with a separated soul's unhindered entrance into beatitude\\
\textbf{6.} Whether the Angels Receive Grace and Glory According to the Degree of Their Natural Gifts? & Aquinas reasonably proposes grace and glory proportioned to angelic natural gifts. & Thomist position. Peter Lombard, Sentences II d.3\\
\textbf{7.} Whether Natural Knowledge and Love Remain in the Beatified Angels? & Natural knowledge and love remain in beatitude, ordered to higher knowledge and love. & Thomist position. Perfection preserves nature and its operations\\
\textbf{8.} Whether a Beatified Angel Can Sin? & The beatified angel cannot sin, because it sees God as goodness itself. & Common teaching. Augustine, Literal Meaning of Genesis XI\\
\textbf{9.} Whether the Beatified Angels Advance in Beatitude? & Essential beatitude does not increase; joy concerning the salvation of others can increase until judgment. & Thomist position. Merit belongs to wayfarers, not possessors of the end\\
\bottomrule
\end{longtable}
\Needspace{10\baselineskip}
\subsection{Question 63}
\begin{longtable}{@{}>{\raggedright\arraybackslash}p{0.31\linewidth}>{\raggedright\arraybackslash}p{0.39\linewidth}>{\raggedright\arraybackslash}p{0.23\linewidth}@{}}
\toprule
Article and question & Aquinas's determination & Grade; authority\\
\midrule\endhead
\textbf{1.} Whether the Evil of Fault Can Be in the Angels? & A created angel can sin by willing without conformity to God's rule. & Church teaching. Job 4:18\\
\textbf{2.} Whether Only the Sin of Pride and Envy Can Exist in an Angel? & The first angelic sin is pride; envy follows; demons incur guilt in other sins by inducing them. & Common teaching. Augustine, City of God XIV.3\\
\textbf{3.} Whether the Devil Desired to Be As God? & The devil sought self-sufficient final happiness or supernatural likeness by its own power, not literal equality of essence with God. & Thomist position. Isaiah 14:13--14; Questions on Old and New Testament 113, attributed to Augustine in the Summa\\
\textbf{4.} Whether Any Demons Are Naturally Wicked? & No demon is evil by nature; wickedness is voluntary. & Defined. Dionysius, Divine Names IV\\
\textbf{5.} Whether the Devil Was Wicked by the Fault of His Own Will in the First Instant of His Creation? & The devil did not sin in the first instant of creation; the initial operation comes from the good Creator. & Thomist position. Genesis 1:31\\
\textbf{6.} Whether There Was Any Interval Between the Creation and the Fall of the Angel? & Aquinas judges the fall immediately after the first instant more probable, given creation in grace and immediate reward. & Disputed. John 8:44, interpreted through Augustine, City of God XI.15\\
\textbf{7.} Whether the Highest Angel Among Those Who Sinned Was the Highest of All? & Aquinas favors the chief sinner's being the highest angel, while permitting the earthly-ruler view. & Disputed. Gregory the Great, Gospel homily 34\\
\textbf{8.} Whether the Sin of the Highest Angel Was the Cause of the Others Sinning? & The chief angel induced others to sin by exhortation, without coercing their choice. & Thomist position. Apocalypse 12:4\\
\textbf{9.} Whether Those Who Sinned Were As Many As Those Who Remained Firm? & More angels remained faithful than sinned; no exact numerical census is determined. & Thomist position. 4 Kings 6:16\\
\bottomrule
\end{longtable}
\Needspace{10\baselineskip}
\subsection{Question 64}
\begin{longtable}{@{}>{\raggedright\arraybackslash}p{0.31\linewidth}>{\raggedright\arraybackslash}p{0.39\linewidth}>{\raggedright\arraybackslash}p{0.23\linewidth}@{}}
\toprule
Article and question & Aquinas's determination & Grade; authority\\
\midrule\endhead
\textbf{1.} Whether the Demons' Intellect Is Darkened by Privation of the Knowledge of All Truth? & Demonic natural knowledge remains; revealed information is limited; loving supernatural wisdom is absent. & Thomist position. Dionysius, Divine Names IV\\
\textbf{2.} Whether the Will of the Demons Is Obstinate in Evil? & Demons are fixed in their chosen aversion, though still capable of choosing diverse means. & Church teaching. Psalm 73:23 [Hebrew 74:23]\\
\textbf{3.} Whether There Is Sorrow in the Demons? & Demons have sorrow as opposition of will to unwelcome realities, not bodily passion or repentance for fault as fault. & Thomist position. Apocalypse 18:7\\
\textbf{4.} Whether Our Atmosphere Is the Demons' Place of Punishment? & Hell is their due punishment; present permitted activity in the dark air does not suspend punishment. & Thomist position. Augustine, Literal Meaning of Genesis III.10\\
\bottomrule
\end{longtable}
\Needspace{10\baselineskip}
\subsection{Question 106}
\begin{longtable}{@{}>{\raggedright\arraybackslash}p{0.31\linewidth}>{\raggedright\arraybackslash}p{0.39\linewidth}>{\raggedright\arraybackslash}p{0.23\linewidth}@{}}
\toprule
Article and question & Aquinas's determination & Grade; authority\\
\midrule\endhead
\textbf{1.} Whether one angel enlightens another? & One angel enlightens another by strengthening the recipient intellect and proportioning a communicated truth; God alone gives nature, grace, glory, and immediate vision of his essence. & Thomist position. Dionysius, Celestial Hierarchy VIII\\
\textbf{2.} Whether one angel moves another angel's will? & An angel can persuade another through a presented good but cannot inwardly cause the act of willing. & Thomist position. Theological argument: only God bestows righteousness and inwardly changes the will\\
\textbf{3.} Whether an inferior angel can enlighten a superior angel? & Inferior angels do not illuminate superiors; the heavenly order differs from earthly ecclesiastical office. & Thomist position. Dionysius, Celestial Hierarchy IV; Ecclesiastical Hierarchy V\\
\textbf{4.} Whether the superior angel enlightens the inferior as regards all he himself knows? & Superior angels communicate their gifts as far as recipients can receive them; unequal modes of understanding and new disclosures remain. & Thomist position. Gregory through Peter Lombard, Sentences II d.9; Dionysius, Celestial Hierarchy XV\\
\bottomrule
\end{longtable}
\Needspace{10\baselineskip}
\subsection{Question 107}
\begin{longtable}{@{}>{\raggedright\arraybackslash}p{0.31\linewidth}>{\raggedright\arraybackslash}p{0.39\linewidth}>{\raggedright\arraybackslash}p{0.23\linewidth}@{}}
\toprule
Article and question & Aquinas's determination & Grade; authority\\
\midrule\endhead
\textbf{1.} Whether one angel speaks to another? & An angel speaks by voluntarily directing a mental conception to another, without a bodily sign. & Thomist position. 1 Corinthians 13:1\\
\textbf{2.} Whether the inferior angel speaks to the superior? & Inferiors can speak to superiors; communication of a private intention is speech without illumination. & Thomist position. Dionysius, Celestial Hierarchy VII (Psalm 23[24])\\
\textbf{3.} Whether an angel speaks to God? & Angels speak to God in praise, admiration, and consultation, not to supply knowledge he lacks. & Thomist position. Zacharias 1:12\\
\textbf{4.} Whether local distance influences the angelic speech? & Local distance does not obstruct angelic speech, which is an intellectual act. & Thomist position. Luke 16:24\\
\textbf{5.} Whether all the angels know what one speaks to another? & An angel may address one recipient without communicating the same conception to every angel. & Thomist position. Argument from the possibility of private human speech\\
\bottomrule
\end{longtable}
\Needspace{10\baselineskip}
\subsection{Question 108}
\begin{longtable}{@{}>{\raggedright\arraybackslash}p{0.31\linewidth}>{\raggedright\arraybackslash}p{0.39\linewidth}>{\raggedright\arraybackslash}p{0.23\linewidth}@{}}
\toprule
Article and question & Aquinas's determination & Grade; authority\\
\midrule\endhead
\textbf{1.} Whether all the angels are of one hierarchy? & One divine ruler governs all; three angelic hierarchies are distinguished by modes of receiving illumination about divine works. & Thomist position. Dionysius, Celestial Hierarchy VI\\
\textbf{2.} Whether there are several orders in one hierarchy? & Each hierarchy contains three orders distinguished by offices: beginning, middle, and completion. & Thomist position. Ephesians 1:20--21\\
\textbf{3.} Whether there are many angels in one order? & Many angels belong to an order under the general classification available to us; each has a more particular office. & Thomist position. Isaias 6:3\\
\textbf{4.} Whether the distinction of hierarchies and orders comes from the angelic nature? & Orders involve natural gifts and gracious gifts; supernatural fruition formally depends upon grace. & Thomist position. Peter Lombard, Sentences II d.9\\
\textbf{5.} Whether the orders of the angels are properly named? & Scriptural order names signify proper excellence; shared perfections are possessed properly, eminently, or by participation. & Thomist position. Scriptural names: Isaias 6; Ezechiel 1/10; Colossians 1:16; Ephesians 1:21; Jude 9\\
\textbf{6.} Whether the grades of the orders are properly assigned? & Dionysius and Gregory exchange the ranks of Principalities and Virtues; both arrangements have intelligible functional explanations. & Disputed. Dionysius, Celestial Hierarchy VII--IX\\
\textbf{7.} Whether the orders will outlast the Day of Judgment? & Differences of nature and glory endure; missions leading pilgrims to their end cease in that form, while communion and illumination remain. & Thomist position. Judges 5:20, applied to angels\\
\textbf{8.} Whether men are taken up into the angelic orders? & Humans can attain angelic degrees of glory by grace without acquiring angelic nature. & Thomist position. Matthew 22:30\\
\bottomrule
\end{longtable}
\Needspace{10\baselineskip}
\subsection{Question 109}
\begin{longtable}{@{}>{\raggedright\arraybackslash}p{0.31\linewidth}>{\raggedright\arraybackslash}p{0.39\linewidth}>{\raggedright\arraybackslash}p{0.23\linewidth}@{}}
\toprule
Article and question & Aquinas's determination & Grade; authority\\
\midrule\endhead
\textbf{1.} Whether there are orders among the demons? & Demons retain natural order; they lost initial grace and never possessed consummated glory. & Thomist position. Ephesians 6:12\\
\textbf{2.} Whether among the demons there is precedence? & Demonic precedence follows remaining natural inequalities; cooperation arises from shared malice rather than charity. & Thomist position. Gloss on 1 Corinthians 15:24\\
\textbf{3.} Whether there is enlightenment in the demons? & Demons communicate conceptions but do not illuminate one another by ordering truth to God. & Thomist position. Ecclesiasticus 34:4 and the account of illumination\\
\textbf{4.} Whether the good angels have precedence over the bad angels? & Good angels govern bad angels by participation in divine justice, even where a demon has greater natural powers. & Thomist position. Augustine, De Trinitate III.4.9 (received citation says II); Gregory, Homilies II.34.10\\
\bottomrule
\end{longtable}
\Needspace{10\baselineskip}
\subsection{Question 110}
\begin{longtable}{@{}>{\raggedright\arraybackslash}p{0.31\linewidth}>{\raggedright\arraybackslash}p{0.39\linewidth}>{\raggedright\arraybackslash}p{0.23\linewidth}@{}}
\toprule
Article and question & Aquinas's determination & Grade; authority\\
\midrule\endhead
\textbf{1.} Whether the corporeal creature is governed by the angels? & Bodily creatures are governed through angels within providence; distinctions of office do not require a choir per bodily species. & Thomist position. Augustine, De Trinitate III.4; Gregory, Dialogues IV.6\\
\textbf{2.} Whether corporeal matter obeys the mere will of an angel? & Matter does not take any substantial form by an angel's mere willing; angels arrange bodily agents. & Thomist position. Augustine, De Trinitate III.8.13 (locus supplied from direct reading)\\
\textbf{3.} Whether bodies obey the angels as regards local motion? & Angels can move bodies locally and thereby bring other bodily changes about through natural causes. & Thomist position. Augustine, De Trinitate III.8--9\\
\textbf{4.} Whether angels can work miracles? & God alone works miracles beyond the whole created order; angels intercede or minister rather than accomplish them by natural power. & Thomist position. Psalm 135:4 (Hebrew 136:4)\\
\bottomrule
\end{longtable}
\Needspace{10\baselineskip}
\subsection{Question 111}
\begin{longtable}{@{}>{\raggedright\arraybackslash}p{0.31\linewidth}>{\raggedright\arraybackslash}p{0.39\linewidth}>{\raggedright\arraybackslash}p{0.23\linewidth}@{}}
\toprule
Article and question & Aquinas's determination & Grade; authority\\
\midrule\endhead
\textbf{1.} Whether an angel can enlighten man? & Angels strengthen human understanding and propose truth through sensible likenesses; infused faith remains God's gift. & Thomist position. Dionysius, Celestial Hierarchy IV\\
\textbf{2.} Whether the angels can change the will of man? & Angels can persuade and affect passions but cannot inwardly cause or necessitate the human will's assent. & Thomist position. Proverbs 21:1\\
\textbf{3.} Whether an angel can change man's imagination? & Good and bad angels can change imagination through bodily conditions; they use material received through sense. & Thomist position. Matthew 1:20; 2:13,19\\
\textbf{4.} Whether an angel can change the human senses? & Angels can change sensation through external objects or internal bodily conditions. & Thomist position. Genesis 19:11; 4 Kings 6:18\\
\bottomrule
\end{longtable}
\Needspace{10\baselineskip}
\subsection{Question 112}
\begin{longtable}{@{}>{\raggedright\arraybackslash}p{0.31\linewidth}>{\raggedright\arraybackslash}p{0.39\linewidth}>{\raggedright\arraybackslash}p{0.23\linewidth}@{}}
\toprule
Article and question & Aquinas's determination & Grade; authority\\
\midrule\endhead
\textbf{1.} Whether the angels are sent on works of ministry? & Angels are sent by divine command to apply their limited power in ministry, without losing contemplation or dignity. & Church teaching. Exodus 23:20\\
\textbf{2.} Whether all the angels are sent in ministry? & The higher orders are never sent on direct external missions, although all communicate intellectual effects. & Disputed. Gregory, Homilies II.34, reporting Dionysius, Celestial Hierarchy XIII\\
\textbf{3.} Whether all the angels who are sent, assist? & All blessed angels assist by seeing God directly; only the highest hierarchy assists in the narrower sense of receiving divine mysteries immediately. & Thomist position. Gregory, Moralia XVII on Job 25:3\\
\textbf{4.} Whether all the angels of the second hierarchy are sent? & Five orders execute external ministry; Dominations direct it, while the first three orders and Dominations do not execute external missions. & Thomist position. Dionysius, Celestial Hierarchy VIII\\
\bottomrule
\end{longtable}
\Needspace{10\baselineskip}
\subsection{Question 113}
\begin{longtable}{@{}>{\raggedright\arraybackslash}p{0.31\linewidth}>{\raggedright\arraybackslash}p{0.39\linewidth}>{\raggedright\arraybackslash}p{0.23\linewidth}@{}}
\toprule
Article and question & Aquinas's determination & Grade; authority\\
\midrule\endhead
\textbf{1.} Whether men are guarded by the angels? & Angels guard fallible human persons by helping practical understanding; their help neither replaces grace nor forces consent. & Church teaching. Psalm 90:11 (Hebrew 91:11)\\
\textbf{2.} Whether each man is guarded by an angel? & An individual guardian is assigned to each human person because providence concerns the enduring individual rational soul. & Thomist position. Jerome, Commentary on Matthew III, at 18:10, cited by Aquinas\\
\textbf{3.} Whether to guard men belongs only to the lowest order of angels? & Personal guardians belong to the lowest order; broader guardianships extend through higher orders. & Thomist position. Psalm 90 (Hebrew 91); Dionysius, Celestial Hierarchy V, IX\\
\textbf{4.} Whether angels are appointed to the guardianship of all men? & Every human wayfarer receives angelic protection, including unbelievers; Christ has ministering angels rather than a superior guardian. & Church teaching. Jerome, Commentary on Matthew III, at 18:10, cited by Aquinas\\
\textbf{5.} Whether an angel is appointed to guard a man from his birth? & Personal guardianship begins at birth rather than baptism; maternal guardianship of the unborn is offered as probable. & Thomist position. Jerome, Commentary on Matthew III, at 18:10, cited by Aquinas\\
\textbf{6.} Whether the angel guardian ever forsakes a man? & A guardian never abandons a person entirely, though providence can permit particular harms or sins. & Thomist position. 1 Peter 5:8, a fortiori argument for unceasing good guardianship\\
\textbf{7.} Whether angels grieve for the ills of those whom they guard? & Blessed angels do not suffer grief; they reject sin while adhering to divine justice and permission. & Thomist position. Apocalypse 21:4\\
\textbf{8.} Whether there can be strife or discord among the angels? & Resistance among good angels concerns opposed matters presented for divine judgment, not contrary wills. & Thomist position. Daniel 10:13\\
\bottomrule
\end{longtable}
\Needspace{10\baselineskip}
\subsection{Question 114}
\begin{longtable}{@{}>{\raggedright\arraybackslash}p{0.31\linewidth}>{\raggedright\arraybackslash}p{0.39\linewidth}>{\raggedright\arraybackslash}p{0.23\linewidth}@{}}
\toprule
Article and question & Aquinas's determination & Grade; authority\\
\midrule\endhead
\textbf{1.} Whether men are assailed by the demons? & Demonic malice causes the assault; its permission and ordering toward good belong to divine providence. & Church teaching. Ephesians 6:12\\
\textbf{2.} Whether to tempt is proper to the devil? & To tempt is to test; demons test to harm and infer inward dispositions from effects rather than knowing hearts directly. & Thomist position. 1 Thessalonians 3:5 and its gloss\\
\textbf{3.} Whether all sins are due to the temptation of the devil? & Not every sin is directly instigated by a demon; the devil is an indirect cause through the first temptation and its consequences. & Thomist position. De ecclesiasticis dogmatibus 49 (received numbering)\\
\textbf{4.} Whether demons can lead men astray by means of real miracles? & Demons cannot work miracles strictly understood; real natural effects and deceptive appearances can nevertheless mislead. & Thomist position. Augustinian attribution; English edition marks Liber XXI sententiarum 4 as supposititious\\
\textbf{5.} Whether a demon who is overcome by man, is for this reason hindered from making further assaults? & Temporary withdrawal after defeat is more probable than permanent exclusion; the tempter may return with permission. & Disputed. Matthew 4:11\\
\bottomrule
\end{longtable}
